\documentclass[conference]{IEEEtran}
\IEEEoverridecommandlockouts

% ---- packages ----
\usepackage{cite}
\usepackage{amsmath,amssymb,amsfonts}
\usepackage{algorithmic}
\usepackage{graphicx}
\usepackage{booktabs}
\usepackage{array}
\usepackage{textcomp}
\usepackage{xcolor}
\usepackage{url}
\usepackage[hidelinks]{hyperref}

% convenient macros
\newcommand{\code}[1]{\texttt{#1}}

\begin{document}

\title{Reproduction and Architectural Ablation of PPO and DQN for
Discrete-Action Autonomous Highway Driving}

\author{\IEEEauthorblockN{Shao-Ming Ruan}
\IEEEauthorblockA{\textit{Institute of Electrical Engineering (Group F)} \\
\textit{National Cheng Kung University}\\
Tainan, Taiwan}
}

\maketitle

\begin{abstract}
Decision-making for autonomous driving is naturally framed as a sequential
decision problem, making it a strong testbed for deep reinforcement learning
(RL). This project \emph{reproduces and ablates} two canonical deep-RL
algorithms---Proximal Policy Optimization (PPO) and Double Deep Q-Network
(Double DQN)---on the discrete-action \code{highway-env} driving simulator. The
goal is to investigate two practical design choices: \emph{(i)} does a
permutation-aware Self-Attention feature extractor outperform a plain
Multi-Layer Perceptron (MLP) under a strict parameter-matched budget, and
\emph{(ii)} how do on-policy and off-policy methods compare in both sample
efficiency and resulting driving behaviour?

Results indicate that under the aggressive reward style ($n=5$ seeds), PPO with
Self-Attention ($235.7 \pm 16.7$), PPO with MLP ($223.0 \pm 19.6$), and Double
DQN ($234.1 \pm 7.7$) all approach a comparable theoretical reward ceiling
(${\approx}266.7$). However, Double DQN achieves the lowest cross-seed variance
($\sigma=7.7$) and faster initial learning. \textbf{Critically, behavioural
evaluation reveals a stark contrast}: Double DQN achieves a significantly lower
crash rate ($11.3\%$) compared to the extreme bimodal behaviours observed in PPO
variants ($60.7$--$80.0\%$ crash rates). Because these metrics are measured from
deterministic policies and the reward styles differ only in reward coefficients,
this gap reflects a genuine difference in the \emph{learned policies}, not a
test artifact---demonstrating that episodic return alone is an insufficient
metric for safety-critical environments. This pattern \emph{generalises} to a
second environment (\code{merge-v0}), where the three methods tie even more
tightly on return ($\sigma\leq0.3$) yet PPO crashes in $100\%$ of episodes
against DQN's $0\%$. A
reward-shaping study further confirms that qualitative driving style is
overwhelmingly determined by reward design rather than architectural choice, yet
convergence to safe behavioural basins remains highly sensitive to optimisation
variance. All code and trained models are released for full reproducibility.
\end{abstract}

\begin{IEEEkeywords}
Reinforcement learning, Proximal Policy Optimization, Deep Q-Network,
self-attention, autonomous driving, ablation study, highway-env.
\end{IEEEkeywords}

\section{Introduction}
Autonomous-driving decision-making---deciding \emph{when} to change lanes,
accelerate, or yield---can be modelled as a Markov Decision Process (MDP) in
which an agent repeatedly observes nearby traffic and selects a high-level
manoeuvre. Deep reinforcement learning (RL) is an attractive solution because
it learns a closed-loop policy directly from interaction, without
hand-engineering the trade-off between progress and safety. The
\code{highway-env} simulator has become a standard, lightweight benchmark for
exactly this class of problems: it exposes a small kinematic state, a discrete
set of meta-actions, and a configurable reward that trades speed against
collisions.

A practitioner approaching this benchmark faces two recurring design choices
that are rarely studied in a controlled way:
\begin{enumerate}
\item \textbf{Feature extractor.} The observation is a \emph{set} of vehicles
(the main vehicle and its $N$ nearest neighbours). A natural inductive bias is
to process this set with \textbf{self-attention}, which is
permutation-equivariant and lets the main vehicle attend to the most relevant
neighbours. The simpler baseline is to \textbf{flatten} the set and feed a
plain MLP. Does the attention bias actually help when the two networks are given
the \emph{same} parameter budget and the \emph{same} training pipeline?
\item \textbf{Algorithm family.} PPO (on-policy) and DQN (off-policy) are the
two workhorses of deep RL. Their relative sample efficiency and final behaviour
on this benchmark are folklore; we measure them under identical conditions.
\end{enumerate}

\textbf{Positioning.} Applying self-attention to an ego-plus-$N$-vehicle
observation is a \emph{standard} design in \code{highway-env}; reproducing it is
not a claim to novelty. The contribution is an \textbf{empirical finding with
methodological teeth}: under a controlled, parameter-fair comparison we show that
\emph{episodic return---the metric almost universally reported on this
benchmark---can be actively misleading for safety-critical control}, and we
quantify by how much.

\textbf{Contributions.}
\begin{itemize}
\item \textbf{The central finding.} Across a parameter-fair ablation, all three
configurations land within $5.7\%$ on episodic return, yet their deployed crash
rates differ by an order of magnitude ($11.3\%$ for DQN vs.\ $60$--$80\%$ for
PPO). This is replicated on a \emph{second} environment (\code{merge-v0}), where
returns tie within $1.4\%$ but PPO crashes in $100\%$ of episodes against DQN's
$0\%$. We argue this makes \textbf{behaviour-level evaluation a necessary
complement to return curves} on this benchmark, not an optional extra.
\item A \textbf{parameter-matched architectural ablation}---Self-Attention
vs.\ MLP actor--critic ($66.8$k vs.\ $72.3$k parameters, a $7.5\%$ gap)---isolating
the inductive bias from model capacity, alongside an off-policy Double-DQN
baseline.
\item A from-scratch, readable implementation of \textbf{PPO} (GAE, clipped
surrogate, per-minibatch advantage normalisation, linear LR decay, gradient
clipping) and \textbf{Double DQN} (replay buffer, $\epsilon$-greedy, target
network, double-Q target), cross-checked against \textbf{Stable-Baselines3}
(Section~\ref{sec:methods}).
\item \textbf{Multi-seed} training ($n=5$) with \textbf{mean$\pm$std} learning
curves, and a \textbf{reward-shaping study} (conservative / base / aggressive)
showing reward design---not architecture---dominates qualitative driving style.
\end{itemize}

All claims are scoped to the experimental conditions studied.

\section{Background and Related Work}
\textbf{Proximal Policy Optimization (PPO)}~\cite{schulman2017proximal} is an
on-policy policy-gradient method that maximises a \emph{clipped surrogate
objective} to keep each update close to the data-collecting policy, giving much
of the stability of trust-region methods at first-order cost. It is combined
with \textbf{Generalized Advantage Estimation (GAE)}~\cite{schulman2015high},
which interpolates between high-bias/low-variance TD and low-bias/high-variance
Monte-Carlo advantage estimates via parameter $\lambda$.

\textbf{Deep Q-Networks (DQN)}~\cite{mnih2015human} learn an off-policy value
function with experience replay and a target network. \textbf{Double
DQN}~\cite{van2016deep} decouples action \emph{selection} from action
\emph{evaluation} in the bootstrap target to reduce the well-known
over-estimation bias of vanilla Q-learning.

\textbf{Attention for driving.} Leurent and Mercat~\cite{leurent2019social}
proposed a \emph{social attention} architecture for tactical decision-making in
dense traffic, letting the ego vehicle attend over a variable number of
neighbours. This design is now a common baseline shipped with \code{highway-env}.
Here it is reproduced and used as one arm of a controlled ablation.

\textbf{Benchmark.} \code{highway-env}~\cite{leurent2018environment} is a
collection of tactical-driving tasks (\code{highway}, \code{merge},
\code{roundabout}, \code{intersection}) with configurable observation, action,
dynamics, and reward. The \code{Kinematics} observation and the
\code{DiscreteMetaAction} action space are used throughout.

\section{Problem Formulation}
Each task is modeled as an episodic MDP $(S, A, P, r, \gamma)$.

\textbf{State / observation.} The \code{Kinematics} observation is a matrix of
shape $V \times F$ with $V = 5$ vehicles (the ego vehicle plus its four nearest
neighbours) and $F = 5$ features per vehicle:
\code{presence}, $x$, $y$, $v_x$, $v_y$. Coordinates and velocities are
ego-relative (\code{absolute=False}) and normalised to roughly $[-1, 1]$
(\code{normalize=True}). The observation is therefore a small \emph{set} of
vehicles, motivating the permutation-aware attention model.

\textbf{Action.} \code{DiscreteMetaAction} exposes five high-level manoeuvres:
$0=$\,\code{LANE\_LEFT}, $1=$\,\code{IDLE}, $2=$\,\code{LANE\_RIGHT},
$3=$\,\code{FASTER}, $4=$\,\code{SLOWER}. The agent acts at
\code{policy\_frequency}~$=5$~Hz while physics runs at
\code{simulation\_frequency}~$=15$~Hz.

\textbf{Episode horizon (a critical environment characteristic).}
\code{highway-env}'s \code{duration} is measured in \emph{simulated seconds},
not steps. Truncation fires when \code{self.time}~$\geq$~\code{duration}, and
\code{self.time} advances by $1/\text{\code{policy\_frequency}} = 0.2$~s per
decision. Therefore \code{duration}~$=80$ yields a \textbf{maximum of
$80 \times 5 = 400$ decision steps per episode}, not 80 steps. Training executes
$100{,}000$ environment steps $\approx 97$ PPO updates.

\textbf{Reward shaping.} \code{highway-env}'s reward trades high speed against
collisions and lane-change penalties. Three \emph{driving styles} are exposed
that change the reward coefficients (Table~\ref{tab:reward}). These coefficients
were tuned empirically rather than chosen analytically: with the high-speed
reward raised to the aggressive value of $2.0$ but the collision penalty kept at
the default $-1.0$, agents largely ignored collisions and maintained maximum
speed; raising the penalty to $-4.0$ was required before collision avoidance
registered in the learned policy.
Crucially, the \textbf{absolute return scale differs
between styles} (e.g.\ aggressive multiplies the high-speed term), so episodic
returns are only comparable \emph{within} a style. Cross-style comparison is
performed with behaviour metrics instead.

\begin{table}[t]
\caption{Reward coefficients per driving style (\code{highway-v0}).}
\label{tab:reward}
\centering
\small
\begin{tabular}{@{}lrrrrr@{}}
\toprule
Style & coll. & hi-spd & lc & rt-lane & spd (m/s) \\
\midrule
base         & $-1.0$ & $0.4$ & $0.0$  & def. & $[20,30]$ \\
conservative & $-2.0$ & $0.4$ & $-0.5$ & def. & $[10,20]$ \\
aggressive   & $-4.0$ & $2.0$ & $+0.2$ & $0.0$ & $[30,40]$ \\
\bottomrule
\end{tabular}
\end{table}

\section{Methods}
\label{sec:methods}

\subsection{PPO (from scratch)}
At each iteration the agent collects a rollout of $T = 1024$ transitions, then
performs $K = 10$ epochs of minibatch SGD. Advantages are computed with GAE:
\begin{align}
\delta_t &= r_t + \gamma\, V(s_{t+1})\,(1 - d_{t+1}) - V(s_t), \\
A_t &= \delta_t + \gamma\lambda\,(1 - d_{t+1})\, A_{t+1}, \\
R_t &= A_t + V(s_t) \quad \text{(value targets)},
\end{align}
with $\gamma = 0.99$, $\lambda = 0.95$. The policy is optimised with the
clipped surrogate objective
\begin{equation}
L^{\mathrm{CLIP}} = \mathbb{E}\!\left[ \min\!\big( \rho_t \hat{A}_t,\;
\mathrm{clip}(\rho_t, 1-\epsilon, 1+\epsilon)\, \hat{A}_t \big) \right],
\end{equation}
\begin{equation}
\rho_t = \exp\!\big( \log \pi_\theta(a_t|s_t) - \log \pi_{\theta_{\mathrm{old}}}(a_t|s_t) \big), \quad \epsilon = 0.2.
\end{equation}
The total loss is
\begin{equation}
L = L^{\mathrm{CLIP}}_{\mathrm{pg}} - c_{\mathrm{ent}}\, H[\pi_\theta]
+ c_{\mathrm{vf}}\, \big( V_\theta(s_t) - R_t \big)^2,
\end{equation}
with entropy coefficient $c_{\mathrm{ent}} = 0.005$ and value coefficient
$c_{\mathrm{vf}} = 0.5$. Gradients are clipped to global norm $0.5$; the
optimiser is Adam ($\text{lr} = 5\times10^{-4}$, $\epsilon = 10^{-5}$) with
learning rate \textbf{linearly decayed} to zero. A critical implementation
detail for stability is handling GAE termination correctly: the per-step
\code{nextnonterminal} mask must be applied so that advantage estimates do not
leak across episode boundaries, and advantage normalisation must be performed
\emph{inside} each minibatch rather than over the whole rollout. Omitting either
produces erratic learning curves.

\subsection{Double DQN (baseline)}
The off-policy baseline stores transitions in a replay buffer ($|B| = 50{,}000$)
and minimises the TD error
\begin{align}
y &= r + \gamma\, Q_{\mathrm{target}}\!\big(s', \arg\max_a Q_{\mathrm{online}}(s', a)\big)(1 - d), \\
L &= \mathrm{MSE}\big( Q_{\mathrm{online}}(s, a),\, y \big).
\end{align}
The target network is hard-updated every $500$ gradient steps; exploration uses
$\epsilon$-greedy with $\epsilon$ decayed multiplicatively ($0.995$ per update)
from $1.0$ to a floor of $0.05$; Adam uses $\text{lr} = 10^{-4}$, batch size
$256$, $\gamma = 0.99$.

\subsection{Network Architectures}
All three networks consume the $5 \times 5$ observation.
\begin{itemize}
\item \textbf{AttentionActorCritic.} A linear projection lifts each vehicle's
$5$ features to a $64$-d embedding; a \code{MultiheadAttention} layer ($4$
heads, \code{batch\_first}) lets vehicles attend to one another; the result is
flattened ($5 \times 64 = 320$) and fed to separate two-hidden-layer
($64$--$64$, \code{tanh}) actor and critic heads.
\item \textbf{MlpActorCritic.} The observation is flattened ($25$-d) and passed
through a $25 \to 64 \to 320$ ReLU trunk, then the same $64$--$64$ \code{tanh}
actor/critic heads. Trunk width is chosen so the total parameter count matches
the attention model.
\item \textbf{MlpQNetwork (DQN).} A compact $25 \to 128 \to 128 \to 5$ ReLU
network.
\end{itemize}

\begin{table}[t]
\caption{Parameter budget.}
\label{tab:params}
\centering
\small
\begin{tabular}{@{}lcc@{}}
\toprule
Network & Parameters & Role \\
\midrule
AttentionActorCritic & $66{,}822$ & PPO actor--critic (attention) \\
MlpActorCritic       & $72{,}262$ & PPO actor--critic (MLP) \\
MlpQNetwork          & $20{,}485$ & DQN value net \\
\bottomrule
\end{tabular}
\end{table}

The attention model has \textbf{$7.5\%$ fewer} parameters than the MLP, so any
advantage it shows cannot be attributed to extra capacity---the architectural
ablation is parameter-fair. The scope of this fairness should be stressed: it
controls capacity \emph{within} the PPO family (Attention vs.\ MLP). The DQN
Q-network is $3.5\times$ smaller and is \emph{not} capacity-matched to the PPO
networks; the PPO-vs-DQN comparison is therefore an \emph{algorithm-family}
comparison, not a controlled-capacity one. DQN is kept small deliberately, as a
lightweight off-policy baseline, and its competitiveness is reported with that
caveat in mind.

\section{Experimental Setup}
\label{sec:setup}
\textbf{Training.} $100{,}000$ environment steps per run; PPO rollout $1{,}024$;
common initial $\text{lr} = 5\times10^{-4}$. Core comparison configurations
(all \code{highway-v0}, aggressive style, \code{dur=80}):
\textbf{PPO+Attention}, \textbf{PPO+MLP}, \textbf{DQN+MLP}. Each is trained
with seeds $\{0, 1, 2, 3, 4\}$ ($n=5$). Python, NumPy, PyTorch (with
\code{cudnn.deterministic=True}) and the environment reset are all seeded.

\textbf{Reproducibility cross-check.} We additionally train \textbf{SB3
PPO}~\cite{sb3} (\code{MlpPolicy}) with hyperparameters matched to our
implementation ($\text{\code{n\_steps}}=1{,}024$, batch$=256$,
$\text{\code{n\_epochs}}=10$, $\gamma=0.99$, $\lambda=0.95$, clip$=0.2$,
$c_{\mathrm{ent}}=0.005$, $c_{\mathrm{vf}}=0.5$,
$\text{\code{max\_grad\_norm}}=0.5$, $\text{lr}=5\times10^{-4}$), seed 0.

\textbf{Evaluation protocol.} Because training returns are nearly tied across
methods, we compare \emph{behaviour}. For each model we run \textbf{30
deterministic} (arg-max) episodes and report \textbf{Crash rate} (\%),
\textbf{Average survival} (steps), \textbf{Average speed} (m/s), and
\textbf{Average lane changes} per episode. \emph{A clarification on what these
metrics depend on:} for a deterministic policy, the action chosen at each state
is fixed by the network, so the resulting trajectory---and hence all four
behaviour metrics---is determined solely by the policy and the environment
dynamics. It was verified in the \code{highway-env} source that the driving
\emph{styles} differ \emph{only} in reward coefficients
(\code{collision\_reward}, \code{high\_speed\_reward}, \code{lane\_change\_reward},
\code{reward\_speed\_range}; Table~\ref{tab:reward}) and \emph{not} in the
observation or the \code{DiscreteMetaAction} target speeds. Consequently these
metrics are \textbf{invariant to the test-time reward style}: a given policy
produces identical crash/survival/speed/lane-change statistics whether scored
under the base or aggressive coefficients. They are reported under the base
configuration purely for concreteness; they measure the \emph{learned policy
itself}, not a test distribution.

\textbf{Hardware / software.} NVIDIA GeForce RTX 4050 Laptop GPU; PyTorch
\code{2.x+cu126}; Gymnasium + \code{highway-env}; Stable-Baselines3
\code{2.2.1}.

\section{Results}
\label{sec:results}

\subsection{Learning curves (core ablation, aggressive style)}
Figure~\ref{fig:curves} shows mean$\pm$std rolling-average episodic reward
across all five seeds for the three core configurations.
Table~\ref{tab:core} summarises the final performance (mean of the last $10\%$
of logged updates) and a sample-efficiency proxy (first environment step at
which the rolling-average reward reaches $150$).

\begin{figure}[t]
\centering
\includegraphics[width=\columnwidth]{logs/comparison_aggressive.png}
\caption{Mean$\pm$std rolling-average episodic reward (seeds 0--4) for the
three core configurations (aggressive style, \code{highway-v0}, $n=5$ seeds).
Shaded bands indicate $\pm$1 std.}
\label{fig:curves}
\end{figure}

\begin{table}[t]
\caption{Core ablation: aggressive style, \code{highway-v0}, $n=5$ seeds (mean$\pm\sigma$, population std, ddof$=0$). Final reward is the mean of the last $10\%$ of logged updates. Median steps to rolling-avg$\geq$150 across 5 seeds.}
\label{tab:core}
\centering
\small
\begin{tabular}{@{}lccccc@{}}
\toprule
Configuration & $n$ & Final reward & Peak & Steps$\to$avg$\geq$150 \\
\midrule
PPO + Attention          & 5 & $235.7 \pm 16.7$          & $266.7$ & ${\approx}10{,}240$ \\
PPO + MLP                & 5 & $223.0 \pm 19.6$          & $266.7$ & ${\approx}9{,}216$ \\
\textbf{DQN + MLP}       & 5 & $\mathbf{234.1 \pm 7.7}$  & $266.7$ & ${\approx}7{,}721$ \\
SB3 PPO (ref.)           & 1 & $260.5$                   & $266.7$ & $8{,}192$ \\
\bottomrule
\end{tabular}
\vspace{2pt}
\par\footnotesize Max achievable reward $\approx266.7$ (this config). DQN is bolded for lowest cross-seed variance.
\end{table}

\textbf{Overall convergence.} The three configurations differ by at most
$5.7\%$ in final mean reward (the widest gap, PPO+MLP vs.\ PPO+Attention, is
$12.7$ points, or $5.7\%$ of PPO+MLP's $223.0$), and each
approaches the approximate theoretical ceiling ($266.7$). This near-parity
indicates that for this task, architecture or algorithm family does not
determine how high the agent can eventually perform; all three discover
comparably high-quality policies. There is a mechanical reason the returns
collapse together: the aggressive style heavily rewards sustained high speed
($+2.0$ per step), so a fast policy that crashes midway can accumulate a return
comparable to a slower policy that survives the full horizon---the single
$-4.0$ terminal penalty is small relative to the hundreds of steps of
accumulated speed reward. (A crash does truncate the episode, so an \emph{early}
crash, which forfeits that reward, remains costly; it is \emph{mid-to-late}
crashes that the return fails to separate from safe driving.) This is the first
hint of why return and behaviour decouple so sharply in Section~\ref{sec:behav}.

\textbf{Variance across seeds.} The most striking finding is the
\emph{ordering of cross-seed stability}: Double DQN achieves the lowest
variance ($\sigma=7.7$), followed by PPO+Attention ($\sigma=16.7$) and
PPO+MLP ($\sigma=19.6$). DQN's experience-replay buffer smooths the high
variance of early collision-heavy exploration, producing more consistent
convergence trajectories across random seeds. PPO+Attention's modestly lower
variance than PPO+MLP ($16.7$ vs.\ $19.6$) may reflect the permutation-equivariant
inductive bias providing marginally more consistent gradients, but the
difference is small and any strong claim would require more seeds and a
formal significance test.

\textbf{PPO+Attention vs.\ PPO+MLP.} The mean reward difference ($235.7$
vs.\ $223.0$, a $5.7\%$ gap) indicates PPO+Attention reaches higher average
performance. However, per-seed results reveal high variability: PPO+Attention
seed 2 achieves $207.0$ (well below its mean) while seed 1 achieves $255.3$
(well above), and similar spread exists for PPO+MLP. This bimodal outcome---some
seeds converging to effective policies, others not---underlines that the modest
aggregate advantage of attention should not be treated as robust without
$\geq$5 seeds and hypothesis testing.

\textbf{DQN vs.\ PPO.} Double DQN achieves a final mean reward of $234.1$,
placing it \emph{on par} with PPO+Attention ($235.7$). This is notable given
that DQN uses a $3.5\times$ smaller network ($20{,}485$ parameters vs.\
${\sim}67$--$72$k for PPO). Critically, DQN reaches a rolling-avg reward of
$150$ at a \textbf{median of ${\approx}7{,}721$} environment steps---roughly
\textbf{25\% fewer steps} than PPO+Attention's ${\approx}10{,}240$ (and
\textbf{16\%} fewer than PPO+MLP's ${\approx}9{,}216$)---indicating
faster initial learning. The per-update granularity differs (PPO logs at
1,024-step rollout boundaries while DQN logs at episode boundaries), but the
total environment-step count is measured consistently.

\subsection{Behaviour-level evaluation}
\label{sec:behav}
Table~\ref{tab:eval} reports four behaviour metrics over 30 deterministic
episodes per model. All models were trained under the aggressive style. As noted
in Section~\ref{sec:setup}, these metrics reflect the learned policy itself and
are invariant to the test-time reward style; they therefore expose how each
policy actually \emph{drives}.

\begin{table}[t]
\caption{Objective behaviour evaluation (30 episodes, base test config, \code{dur=80}, max 400 steps). Results are averaged across seeds 0--4 ($n=5$).}
\label{tab:eval}
\centering
\small
\begin{tabular}{@{}lcccc@{}}
\toprule
Model & Crash (\%) & Survival & Spd (m/s) & Lane chg. \\
\midrule
PPO + Attention & $60.7$        & $199.9$ & $22.99$ & $25.23$ \\
PPO + MLP       & $80.0$        & $152.6$ & $23.98$ & $16.00$ \\
\textbf{DQN + MLP} & $\mathbf{11.3}$ & $368.5$ & $20.89$ & $155.15$ \\
\bottomrule
\end{tabular}
\vspace{2pt}
\parbox{\columnwidth}{\footnotesize Bold = lowest crash rate. PPO models exhibit bimodal behaviour:
passive-safe in some seeds ($0$--$3.3\%$ crash, $0$ lane chg.) vs.\ catastrophically aggressive in
others ($100\%$ crash). DQN is markedly more consistent ($11.3\%$ avg).}
\end{table}

The behaviour metrics reveal a striking divergence despite near-identical training
returns. \textbf{DQN+MLP achieves the lowest crash rate} ($11.3\%$) with
$368.5$ average survival steps. This figure warrants caution. Qualitative
inspection of rendered episodes shows the DQN agent jittering left--right almost
every step rather than smoothly steering around obstacles: it changes lanes
\emph{far} more than any other model ($155$ per episode on average, with two
seeds exceeding $260$---near-continuous lane switching in a $400$-step episode).
Because the aggressive \emph{training} reward \emph{positively} rewards lane
changes ($+0.2$; Table~\ref{tab:reward}), the policy is actively incentivised to
switch lanes, so the low crash rate may be partly an artifact of a degenerate
``swerve-to-survive'' strategy rather than a genuinely intelligent avoidance
policy. The rendered video is released so that this can be judged directly rather
than taking the crash number at face value. In contrast, \textbf{PPO+Attention and PPO+MLP exhibit
extreme bimodal behaviour}: some seeds converge to a fully passive policy
($0$--$3.3\%$ crash, $\approx 0$ lane changes), while others converge to a
catastrophically aggressive policy ($100\%$ crash), yielding high aggregate
crash rates of $60.7\%$ and $80.0\%$ respectively. This bimodal collapse---present
in the majority of PPO seeds evaluated under the base config after aggressive
training---is consistent with PPO's on-policy updates polarising toward local
optima rather than learning a robust mixed strategy.

\subsection{Stable-Baselines3 cross-validation}
This cross-check was motivated by a concrete observation: PPO+Attention seed 2
came in at $207.0$, well below the other seeds, and it was not immediately clear
whether this reflected ordinary seed variance or an error in the hand-written
PPO. Running a trusted reference implementation under matched hyperparameters is
a cheap way to discriminate the two. Because SB3 uses an \code{MlpPolicy}
$[64,64]$ trunk, its closest architectural counterpart is the custom
\textbf{PPO+MLP} ($223.0 \pm 19.6$), not the attention model. The SB3 run
(seed 0, hyperparameters matched to the present implementation) reached a final
rolling-average reward of \textbf{260.5} and a peak of \textbf{266.7}. This
single-seed value exceeds \emph{every} one of our custom PPO seeds (best
individual: PPO+Attention seed 1, $255.3$; best PPO+MLP seed, $248.4$) and sits
${\approx}17\%$ above the PPO+MLP 5-seed mean. This gap is read not as evidence
that the custom PPO is broken---it reaches the same peak ($266.7$) and the same
qualitative learning dynamics (rapid rise before ${\approx}10$k steps, then a
high plateau)---but as a reflection of SB3's mature engineering defaults
(orthogonal initialisation, running observation normalisation, tuned optimiser
$\epsilon$) that were not all replicated. Because SB3 is a single seed against
five, this is a directional sanity check on correctness, not a like-for-like
performance claim; a fully fair comparison would require multi-seed SB3 runs
with the identical feature extractor.

\subsection{Reward-shaping study}
\label{sec:styles}
Beyond the core ablation, PPO+Attention was trained under the \emph{conservative}
and \emph{base} reward styles (seeds 0, 1, and 2; $n=3$) to illustrate how the
\textbf{reward specification dominates qualitative behaviour}.
Table~\ref{tab:styles} shows the final training reward per style.

\begin{table}[t]
\caption{PPO+Attention final training reward by driving style ($n=3$ seeds, mean$\pm$std). \emph{Returns are not comparable across styles} (reward coefficients differ substantially; see Table~\ref{tab:reward}).}
\label{tab:styles}
\centering
\small
\begin{tabular}{@{}lccc@{}}
\toprule
Style & Seeds & Final reward (mean$\pm$std) & Peak \\
\midrule
conservative & 0, 1, 2 & $366.0 \pm 11.0$ & $390.1$ \\
base         & 0, 1, 2 & $265.2 \pm 14.7$ & $297.1$ \\
aggressive   & 0--4    & $235.7 \pm 16.7$ & $266.7$ \\
\bottomrule
\end{tabular}
\end{table}

The apparent ordering (conservative outscoring base and aggressive in absolute
reward) does \textbf{not} imply conservative is the ``best'' policy. The reward
functions differ fundamentally: conservative targets low speed ($[10,20]$~m/s),
penalises lane changes ($-0.5$), and incurs a smaller collision penalty ($-2.0$),
incentivising slow, stable driving that accumulates many small rewards over the
full 400-step horizon. Aggressive targets high speed ($[30,40]$~m/s), rewards
lane changes ($+0.2$), and imposes a heavy collision penalty ($-4.0$); a single
crash can wipe out many steps of reward, depressing absolute return. As
illustrated in Figure~\ref{fig:styles}, the absolute return trajectories
fundamentally decouple early in training---by roughly the first $10\%$ of
updates the conservative curve already sits well above base and aggressive---reflecting
these diverging optimisation targets rather than any difference in algorithmic
capacity, since all three curves are the same PPO+Attention model. The
comparison is \emph{meaningless} across styles; the point of this study is the
\emph{behavioural} contrast revealed by Table~\ref{tab:behavior_styles}.

\begin{table}[t]
\caption{Reward-shaping behavioural comparison (PPO+Attention, 30 episodes, base test config, seeds 0--2 averaged for conservative and base ($n=3$); seeds 0--4 averaged for aggressive ($n=5$)). Higher survival / lower crash indicates safer driving.}
\label{tab:behavior_styles}
\centering
\small
\begin{tabular}{@{}lcccc@{}}
\toprule
Training Style & Crash (\%) & Survival & Spd (m/s) & Lane chg. \\
\midrule
Conservative & $33.3^\dagger$ & $293.3$ & $21.67$ & $26.64$ \\
\textbf{Base} & $\mathbf{0.0}$ & $400.0$ & $20.03$ & $0.0$ \\
Aggressive   & $60.7^\dagger$ & $199.9$ & $22.99$ & $25.23$ \\
\bottomrule
\end{tabular}
\vspace{2pt}
\parbox{\columnwidth}{\footnotesize $^\dagger$Bimodal. Conservative: seeds 0,1 safe (0\% crash, 400
steps), seed 2 catastrophic (100\% crash, 80 steps). Aggressive: seeds 0,3 safe ($\leq$3.3\%),
seeds 1,2,4 catastrophic (100\% crash).}
\end{table}

Even under the same base test configuration, models trained under different
reward styles drive markedly differently (Table~\ref{tab:behavior_styles}): the
base-trained policy is uniformly safe ($0\%$ crash), the aggressive-trained
policy is fastest but crashes most, and the conservative-trained policy splits
bimodally across seeds (the per-seed breakdown is examined in
Sec.~\ref{sec:discussion}). These differences arise from the reward
specification rather than the architecture, confirming reward design as the
primary lever for qualitative driving style.

\begin{figure}[t]
\centering
\includegraphics[width=\columnwidth]{logs/comparison_styles.png}
\caption{PPO+Attention learning curves under three reward styles (conservative,
base, aggressive). Returns are \emph{not} comparable across styles due to
differing reward coefficients (see Table~\ref{tab:reward}).}
\label{fig:styles}
\end{figure}

\subsection{Generalization to a second environment (\code{merge-v0})}
\label{sec:merge}
The \code{merge-v0} environment was selected to test whether the headline
finding---return parity hiding a behavioural gap---is specific to
\code{highway-v0}. Its dynamics are deliberately the \emph{least} like highway:
instead of free-flowing multi-lane cruising, the ego must time a merge from a
short on-ramp onto a highway with an obstacle at the ramp end---a different
traffic topology and failure mode. If the finding survives there, the objection
that ``it only holds on highway'' is difficult to maintain. The core experiment
was repeated with the identical training pipeline and aggressive reward style
($n=3$ seeds per configuration).
Table~\ref{tab:merge} reports the result, and it is \emph{starker} than on
\code{highway-v0}.

\begin{table}[t]
\caption{Generalization to \code{merge-v0} (aggressive style, $n=3$ seeds, train aggressive / test base). Training reward is mean$\pm\sigma$ (ddof$=0$) of the last $10\%$ of updates; behaviour is over 30 deterministic episodes per model under the \code{merge-v0} base config.}
\label{tab:merge}
\centering
\small
\begin{tabular}{@{}lccccc@{}}
\toprule
Config & Train rew. & Crash (\%) & Surv. & Spd & LC \\
\midrule
PPO + Attention & $58.3{\pm}0.2$ & $100.0$ & $34.0$ & $29.94$ & $34.0$ \\
PPO + MLP       & $59.1{\pm}0.3$ & $100.0$ & $35.0$ & $29.58$ & $32.7$ \\
\textbf{DQN + MLP} & $59.0{\pm}0.3$ & $\mathbf{0.0}$ & $77.7$ & $22.07$ & $44.0$ \\
\bottomrule
\end{tabular}
\vspace{2pt}
\parbox{\columnwidth}{\footnotesize All three configurations converge to the same
training reward (within $1.4\%$, $\sigma\leq0.3$), yet at deployment DQN crashes
in $0\%$ of episodes while \emph{both} PPO variants crash in $100\%$---uniformly
across all seeds. \code{merge-v0} episodes end when the ego passes the merge zone
($x>370$) or crashes, so ``Surv.'' reflects reaching the goal (DQN) vs.\ crashing
midway (PPO); crash rate is the cleaner metric here.}
\end{table}

On \code{merge-v0} the three configurations are even more tightly tied on
training reward than on \code{highway-v0}: all converge to ${\approx}58$--$59$
with cross-seed $\sigma\leq0.3$ (versus $\sigma=7.7$--$19.6$ on highway). Note
that \code{merge-v0} normalises its per-step reward to $[0,1]$ and terminates
when the ego passes the ramp ($x>370$) or crashes, so its return scale is not
comparable to highway's and admits no single fixed ceiling; the tight tie simply
means all three reach similar normalised performance. Yet the behavioural split
is total and \emph{deterministic}: both PPO variants converge, in every seed, to
a maximum-speed policy (${\approx}29.9$~m/s) that crashes on $100\%$ of episodes
before completing the merge, whereas Double DQN learns a slower ($22$~m/s) policy
that traverses the ramp without crashing. Unlike highway, there is no bimodality
here---PPO fails uniformly. The same swerve-to-survive caveat applies, and is in
fact sharper here: DQN averages $44$ lane changes over ${\approx}78$ steps (one
every ${\sim}1.8$ steps), and the aggressive training reward pays $+0.2$ per lane
change, so part of DQN's success may be reward-driven lane oscillation rather
than principled avoidance---the supplementary video lets the reader judge. With
that caveat, this second environment \emph{strengthens} the central claim:
near-identical episodic return can co-exist with an all-or-nothing difference in
how the agent drives, and this is not specific to a single benchmark.

\section{Discussion}
\label{sec:discussion}

\textbf{Reward parity vs.\ behavioural divergence.} The core finding of the
ablation is that a parameter-matched attention architecture, a plain MLP, and a
simpler DQN all converge to statistically comparable final episodic returns
($223$--$235$, a maximum gap of $5.7\%$). However, this near-parity in
aggregate reward masks profound differences in the learned policies. As shown in
Table~\ref{tab:eval}, despite similar training returns, DQN achieves a robust
lane-management strategy with an $11.3\%$ crash rate, while PPO variants
collapse into bimodal extremes yielding aggregate crash rates exceeding $60\%$.
This is a concrete, low-dimensional instance of \emph{reward
misspecification}~\cite{amodei2016concrete, pan2022effects}: policies optimise
the scalar return faithfully yet realise qualitatively different---and unsafe---physical
behaviour. It also echoes the reproducibility concerns of Henderson
\emph{et al.}~\cite{henderson2018deep}: aggregate return, reported without
multi-seed behavioural inspection, can paint a misleadingly uniform picture.
\textbf{Episodic return alone is thus an insufficient metric for safety-critical
tasks}---a failure mode that only behaviour-level, multi-seed evaluation exposes.

\textbf{The stabilising effect of off-policy replay.} Contrary to expectations
that a more complex feature extractor (Self-Attention) would inherently yield
more stable control, Double DQN emerges as both the most stable learner
($\sigma=7.7$) and the lowest-crash policy at deployment (subject to the swerving
caveat of Sec.~\ref{sec:behav}). We attribute the \emph{training stability} to the
fundamental difference in how experience is utilised. In the collision-heavy
early learning phase, PPO's on-policy updates are highly susceptible to the
non-stationarity of recent rollouts, leading some seeds to polarise toward
catastrophic local optima (e.g.\ maintaining maximum speed regardless of
obstacles). DQN's experience-replay buffer acts as a stabilising mechanism,
maintaining a diverse distribution of transitions that prevents the Q-network
from overfitting to short-term catastrophic sequences. The attention mechanism
provides only a marginal variance reduction over a plain MLP for PPO
($16.7$ vs.\ $19.6$), suggesting that the off-policy replay mechanism is a far
more dominant factor in policy stability than the inductive bias of the neural
architecture.

\textbf{Sample efficiency: why DQN learns faster initially.} Double DQN reaches
rolling-avg reward $\geq$150 at a median of ${\approx}7{,}721$ environment steps,
roughly \textbf{25\%} earlier than PPO+Attention (${\approx}10{,}240$) and
\textbf{16\%} earlier than PPO+MLP (${\approx}9{,}216$). This aligns with
a well-known property of off-policy learning: experience replay allows DQN to
reuse early collision-heavy transitions many times, smoothing the otherwise
high-variance gradient signal in the early training phase. PPO, which can only
learn from its most-recent rollout, must accumulate sufficient on-policy
experience before gradients stabilise. However, the early-learning advantage of
DQN does not persist to convergence: all methods ultimately reach similar final
rewards (${\sim}223$--$235$).

\textbf{The behavioural gap is intrinsic to the learned policies.} It is worth
being precise about what the crash-rate gap is \emph{not}. Because the policies
are evaluated deterministically and the reward styles differ only in reward
coefficients (not in dynamics or action mapping; Section~\ref{sec:setup}), the
behaviour metrics do not depend on the test-time reward at all---a given network
crashes, or does not, purely as a function of the actions it selects. The
$11.3\%$ vs.\ $60$--$80\%$ gap is therefore not an artifact of a test
distribution; it is a genuine difference in the \emph{policies that were
learned}. PPO (trained under the strong aggressive speed incentive) converges in
most seeds to a maximum-speed policy that physically collides; DQN converges to a
slower policy that does not. We attribute DQN's \emph{training} stability to
experience replay (above); we are careful, however, not to over-claim DQN as
``intrinsically safer'': all models were trained under a single (aggressive)
reward, and the slow, lane-switching policy DQN learned carries its own
swerve-to-survive caveat (Sec.~\ref{sec:behav}). The robust, reward-independent
claim is the headline one: near-tied episodic returns ($223$--$235$) can hide an
order-of-magnitude difference in how the agent actually drives
($11.3\%$ vs.\ $60$--$80\%$ crash).

\textbf{Reward shaping dominates---but variance qualifies the claim.}
Table~\ref{tab:behavior_styles} reveals a nuanced picture. The base-trained
policy (tested under its own training distribution) achieves a crash rate of
only $0.0\%$ and survives $400.0$ steps on average with zero lane changes,
the most consistent and safest profile of all three styles. The
conservative-trained policy, which \emph{intended} to be even safer through a
stricter lane-change penalty ($-0.5$) and lower speed target ($[10,20]$~m/s),
shows a bimodal outcome: seeds 0 and 1 achieve $0\%$ crash (400 steps, no lane
changes, as intended), while seed 2 crashes on every single episode
($100\%$ crash rate, 80 lane changes per episode). The aggregate mean ($33.3\%$
crash) is therefore misleading---it reflects a high variance in policy
convergence rather than an inherently mediocre strategy.

This finding \emph{qualifies} but does not overthrow the reward-shaping
dominance conclusion: reward engineering does constrain the \emph{direction}
of policy learning (conservative configurations tend toward lower-speed, fewer
lane-change optima), but cannot guarantee that all random seeds converge to the
intended behavioural basin. In combination with the high variance already
documented in Table~\ref{tab:core} ($\sigma=16.7$ for PPO+Attention), this
suggests that PPO on this benchmark has a multi-modal loss landscape, and some
seeds get trapped in local optima that are high-reward yet physically unsafe
(high speed, frequent collisions). Reward
engineering remains the \emph{primary} design lever, but practitioners should
validate with multi-seed behavioural evaluation---not just aggregate reward.

\textbf{Implementation correctness.} The custom PPO reaches the same reward
ceiling ($266.7$ peak) and the same qualitative learning dynamics as the SB3
reference, supporting the correctness of the GAE, clipped surrogate, advantage
normalisation, and linear LR decay. We are deliberately cautious here: SB3's
single-seed $260.5$ exceeds all of our PPO seeds, so we claim only that the
implementation is \emph{in the same regime}, not that it matches SB3's tuned
defaults. A definitive correctness check would require multi-seed SB3 runs with
the identical architecture.

\section{Limitations}
\label{sec:limitations}
\begin{itemize}
\item \textbf{Moderate seed count ($n=5$).} Five seeds provide reasonable
variance estimates but are insufficient for rigorous statistical significance
claims; a formal test (e.g.\ Mann-Whitney $U$) requires more runs. Comparative
results are therefore phrased as ``observed under these conditions.''
\item \textbf{Two environments, fewer seeds on the second.} The core ablation is
on \code{highway-v0} ($n=5$); the central finding is corroborated on
\code{merge-v0} ($n=3$, Sec.~\ref{sec:merge}) but not extended to
\code{roundabout} or \code{intersection}.
\item \textbf{Compute- and time-bounded budget.} All runs were done on a single
RTX 4050 laptop GPU, where one $100$k-step training run takes roughly $10$--$22$
minutes wall-clock (the spread depends on how many runs were active concurrently
against limited RAM). That budget is the concrete reason the core comparison
uses $5$ seeds while the \code{merge-v0} generalization study uses only $3$, and
why training is capped at $100$k steps---some configurations may not be fully
converged. Hyperparameters were matched across methods, not individually tuned.
\item \textbf{DQN is intentionally a simplified baseline} (smaller net, no
attention variant, aggressive $\epsilon$-decay); it bounds, rather than
maximises, off-policy performance. It is also \emph{not} capacity-matched to the
PPO networks, so PPO-vs-DQN is an algorithm-family comparison, not a
controlled-capacity one.
\item \textbf{Single training reward for the behavioural comparison.} All models
in the core comparison were trained under the aggressive style; it is not tested
whether DQN would remain the lowest-crash policy if trained under base or
conservative rewards. (Note this is a \emph{training}-side scope limit: the
behaviour metrics themselves are invariant to the test-time reward, so there is
no test-distribution confound; Section~\ref{sec:setup}.)
\item \textbf{Swerve-to-survive ambiguity.} DQN's low crash rate co-occurs with
very frequent lane changes; intelligent avoidance cannot be fully separated from
a degenerate swerving policy without further trajectory analysis (Sec.~\ref{sec:behav}).
\end{itemize}

\section{Conclusion}
This work presented a from-scratch reproduction and parameter-matched architectural
ablation of PPO and Double DQN on the discrete-action \code{highway-env} driving
simulator ($n=5$ seeds per configuration). Three key findings emerge:

\textbf{(1) Return parity across architecture and algorithm.} Under the
aggressive reward style on \code{highway-v0}, all three configurations---PPO
with Self-Attention ($235.7 \pm 16.7$), PPO with MLP ($223.0 \pm 19.6$), and
Double DQN ($234.1 \pm 7.7$)---approach the same theoretical reward ceiling
(${\approx}266.7$), with final means differing by at most $5.7\%$. Architecture
and algorithm family barely affect \emph{how high} the agent scores---which is
exactly what makes episodic return an unreliable basis for choosing between them.

\textbf{(2) Return parity hides an order-of-magnitude behavioural gap.} Despite
near-tied returns, deterministic behavioural evaluation reveals crash rates of
$11.3\%$ (DQN) versus $60.7\%$ (PPO+Attention) and $80.0\%$ (PPO+MLP)---a
difference in the \emph{learned policies} themselves, independent of the
test-time reward. This gap \emph{generalises}: on \code{merge-v0} the three tie
within $1.4\%$ on training return yet PPO crashes in $100\%$ of episodes against
DQN's $0\%$ (Sec.~\ref{sec:merge}). DQN is also the most stable learner
($\sigma=7.7$) and the fastest to reach avg$\geq$150 (median ${\approx}7{,}721$
steps), which we attribute to experience replay smoothing gradient variance. We
stop short of calling DQN ``intrinsically safer'': its low-crash policy is also a
frequent-lane-changer (Sec.~\ref{sec:behav}), and only a single training reward
was tested.

\textbf{(3) Reward shaping as primary lever---with an important variance caveat.}
A conservative-vs.-base-vs.-aggressive reward-shaping study confirms that
reward specification is the dominant determinant of driving style. However,
behavioural evaluation reveals a nuanced finding: even with the same reward
shaping, individual seeds can converge to qualitatively opposite policies (e.g.\
conservative seeds 0,1: $0\%$ crash; seed 2: $100\%$ crash). Reward engineering
constrains the \emph{direction} of learning but does not guarantee all seeds
converge to the intended behavioural basin. Multi-seed behavioural evaluation
is therefore a necessary complement to reward design for safety-critical
deployment.

\textbf{Future work} should \emph{train} each algorithm under several reward
styles (not just aggressive) to test whether DQN's low-crash behaviour is
intrinsic or specific to the aggressive incentive, extend the ablation to
\code{roundabout-v0} and \code{intersection-v0} beyond the \code{highway-v0} and
\code{merge-v0} studied here, and test a prioritised-replay DQN variant to probe
whether the stability advantage strengthens.

\begin{thebibliography}{1}

\bibitem{schulman2017proximal}
J.~Schulman, F.~Wolski, P.~Dhariwal, A.~Radford, and O.~Klimov,
``Proximal policy optimization algorithms,''
\emph{arXiv preprint arXiv:1707.06347}, 2017.

\bibitem{schulman2015high}
J.~Schulman, P.~Moritz, S.~Levine, M.~I. Jordan, and P.~Abbeel,
``High-dimensional continuous control using generalized advantage estimation,''
\emph{arXiv preprint arXiv:1506.02438}, 2015.

\bibitem{mnih2015human}
Mnih, V., Kavukcuoglu, K., Silver, D. \emph{et al.},
``Human-level control through deep reinforcement learning,''
\emph{Nature}, 518, 529–533 (2015). https://doi.org/10.1038/nature14236

\bibitem{van2016deep}
H.~van Hasselt, A.~Guez, and D.~Silver,
``Deep Reinforcement Learning with Double Q-learning,''
in \emph{Proc.\ AAAI}, 2016.

\bibitem{leurent2019social}
E.~Leurent and J.~Mercat,
``Social attention for autonomous decision-making in dense traffic,''
\emph{arXiv preprint arXiv:1911.12250}, 2019.

\bibitem{leurent2018environment}
E.~Leurent,
``An environment for autonomous driving decision-making,''
\emph{GitHub repository}, 2018. https://github.com/Farama-Foundation/HighwayEnv

\bibitem{sb3}
A.~Raffin, A.~Hill, A.~Gleave, A.~Kanervisto, M.~Ernestus, and N.~Dormann,
``Stable-Baselines3: Reliable reinforcement learning implementations,''
\emph{J.\ Mach.\ Learn.\ Res.}, vol.\ 22, no.\ 268, pp.\ 1--8, 2021.

\bibitem{henderson2018deep}
P.~Henderson, R.~Islam, P.~Bachman, J.~Pineau, D.~Precup, and D.~Meger,
``Deep reinforcement learning that matters,''
in \emph{Proc.\ AAAI}, 2018.

\bibitem{amodei2016concrete}
D.~Amodei, C.~Olah, J.~Steinhardt, P.~Christiano, J.~Schulman, and D.~Man\'e,
``Concrete problems in AI safety,''
\emph{arXiv preprint arXiv:1606.06565}, 2016.

\bibitem{pan2022effects}
A.~Pan, K.~Bhatia, and J.~Steinhardt,
``The Effects of Reward Misspecification: Mapping and Mitigating Misaligned Models,''
in \emph{Proc.\ ICLR}, 2022.

\end{thebibliography}

\end{document}
